\Page{Exercices\quad /\quad Découverte}{Sources et modèle}{td1}
\source{Énoncés originaux OptiTAB. Dans tous les exercices, sauf indication contraire, $c=3{,}00\times10^8\ \mathrm{m\,s^{-1}}$ et les lentilles sont minces dans les conditions de Gauss.}
\ex{E01}{Lire une signature spectrale}{Niveau 1}{8 min}
Un spectromètre mesure trois raies à $486{,}1$, $589{,}0$ et $656{,}3\ \mathrm{nm}$ dans le vide.
\begin{questions}
\item Le spectre est-il continu ? La source est-elle monochromatique ?
\item Associer à chaque raie une couleur approximative, puis calculer sa fréquence.
\item Peut-on conclure que la source contient un élément chimique unique à partir de ces seules données ? Justifier.
\end{questions}
\ex{E02}{Une lumière dans le verre}{Niveau 1}{8 min}
Un faisceau de fréquence $5{,}00\times10^{14}\ \mathrm{Hz}$ passe de l'air, assimilé au vide, dans un verre d'indice $1{,}50$.
\begin{questions}
\item Calculer $\lambda_0$, puis la vitesse et la longueur d'onde dans le verre.
\item Quelle grandeur reste inchangée à l'interface immobile ?
\item La longueur d'onde dans le verre est voisine de celle d'une lumière violette dans le vide. Doit-on pour autant dire que le faisceau a changé de couleur spectrale ?
\end{questions}
\ex{E03}{Quand les rayons ne suffisent plus}{Niveau 1}{10 min}
Une onde de longueur d'onde $\lambda=600\ \mathrm{nm}$ éclaire une ouverture de largeur $a$. Pour estimer la diffraction, on fournit $\theta\sim\lambda/a$ ; cette formule n'est qu'un ordre de grandeur.
\begin{questions}
\item Calculer $\lambda/a$ pour $a=1{,}0\ \mathrm{mm}$, puis $a=1{,}0\ \mu\mathrm{m}$.
\item Dans quel cas le modèle des rayons rectilignes au voisinage de l'ouverture est-il le plus adapté ?
\item Pour l'ouverture millimétrique, estimer la demi-largeur de la zone diffractée sur un écran à $2{,}0\ \mathrm{m}$. Expliquer pourquoi « négligeable » dépend de la résolution demandée.
\end{questions}
\begin{methode}{Le réflexe de ce premier groupe}
Écrire chaque conversion en mètres avant de calculer. Une longueur d'onde dans un milieu et une longueur d'onde dans le vide ne se comparent pas directement pour nommer une couleur.
\end{methode}

\Page{Exercices\quad /\quad Découverte}{Traverser une interface}{td2}
\source{Énoncés originaux OptiTAB. Dessiner la normale avant d'écrire une loi.}
\ex{E04}{Le piège des $30^\circ$}{Niveau 1}{10 min}
Un rayon dans l'air arrive sur l'eau ($n=1{,}33$) sous un angle de $30{,}0^\circ$ avec la surface.
\begin{questions}
\item Dessiner la situation et préciser l'angle d'incidence.
\item Calculer l'angle de réfraction et l'angle entre les directions incidente et transmise.
\item Décrire le trajet inverse et préciser l'angle d'incidence correspondant dans l'eau.
\end{questions}
\ex{E05}{Deux faces d'un même verre}{Niveau 1}{12 min}
Un rayon dans un verre d'indice $1{,}52$ arrive à $50{,}0^\circ$ sur une interface, d'abord verre-air, puis verre-eau ($n_e=1{,}33$).
\begin{questions}
\item Calculer les deux incidences critiques et décider s'il existe un rayon transmis dans chaque cas.
\item Quand ce rayon existe, déterminer son angle.
\item Trouver l'intervalle des angles pour lesquels le verre réfléchit totalement au contact de l'air mais transmet au contact de l'eau.
\end{questions}
\ex{E06}{Une lame à faces parallèles}{Niveau 2}{20 min}
Une lame de verre d'épaisseur $e=6{,}00\ \mathrm{mm}$ et d'indice $n=1{,}50$ est entourée d'air. L'incidence est $i=45{,}0^\circ$.
\fig{ex-lame}{Schéma de situation à reproduire en ajoutant les normales et les angles ; dimensions schématiques.}
\begin{questions}
\item Faire un schéma comportant les deux normales et déterminer l'angle interne $r$.
\item Montrer que la direction émergente est parallèle à la direction incidente.
\item Le décalage latéral est la distance perpendiculaire entre ces deux droites. Établir $s=e\sin(i-r)/\cos r$ puis le calculer.
\end{questions}
\Page{Exercices\quad /\quad Découverte}{Images et conventions}{td3}
\ex{E07}{Un miroir qui tourne}{Niveau 2}{15 min}
Un objet est placé à $40{,}0\ \mathrm{cm}$ devant un miroir plan. Un rayon incident fixe se réfléchit sur ce miroir.
\begin{questions}
\item Construire l'image de l'objet et donner sa nature et la distance objet-image.
\item Le miroir pivote de $\beta=3{,}0^\circ$ autour du point d'incidence, dans le plan du dessin. Démontrer que la direction du rayon réfléchi tourne de $2\beta$ en valeur absolue.
\item Expliquer pourquoi une photographie de l'image virtuelle est possible.
\end{questions}
\ex{E08}{Assez net pour quel capteur ?}{Niveau 1}{10 min}
Un point donne une tache de diamètre $12\ \mu\mathrm{m}$. Deux détecteurs ont des pas de pixels respectifs de $5$ et $20\ \mu\mathrm{m}$. On adopte ici le critère simplifié « diamètre inférieur ou égal à un pixel ».
\begin{questions}
\item Pour quel détecteur le point est-il jugé net ? Ce critère signifie-t-il que la tache est physiquement ponctuelle ?
\item Donner les deux conditions de Gauss et leurs conséquences.
\item Pour un rayon de $3{,}0^\circ$, comparer $\tan\theta$ à $\theta$ en radians et calculer l'erreur relative $(\tan\theta-\theta)/\theta$.
\end{questions}
\ex{E09}{Reconnaître une lentille}{Niveau 1}{10 min}
Deux lentilles ont pour focales $f'_1=+8{,}0\ \mathrm{cm}$ et $f'_2=-20{,}0\ \mathrm{cm}$.
\begin{questions}
\item Les classer et calculer leurs vergences en dioptries.
\item Pour chaque lentille, tracer l'axe, $O$, $F$ et $F'$, puis un rayon parallèle à l'axe et un rayon passant par $O$.
\item Expliquer ce que devient un faisceau parallèle incliné et définir le foyer secondaire image correspondant.
\end{questions}
\source{Pour E07 et E09, une construction sans flèche de propagation ou sans distinction entre rayon et prolongement reste incomplète.}

\Page{Exercices\quad /\quad Entraînement}{Trois images à comparer}{td4}
\ex{E10}{Une projection réduite}{Niveau 1}{12 min}
Un objet $AB$ de hauteur $1{,}20\ \mathrm{cm}$ est à $30{,}0\ \mathrm{cm}$ à gauche d'une lentille convergente de focale $10{,}0\ \mathrm{cm}$.
\fig{ex-construction}{Repères de départ pour E10 ; compléter les rayons et l'image. La hauteur est schématique.}
\begin{questions}
\item Donner $\oa$ et calculer $\oap$, $\gamma$ et $\overline{A'B'}$.
\item Caractériser l'image et réaliser une construction à l'échelle.
\end{questions}
\ex{E11}{La même lentille devient une loupe}{Niveau 1}{12 min}
On conserve la lentille de E10, mais l'objet est à $6{,}00\ \mathrm{cm}$ à gauche de $O$.
\begin{questions}
\item Déterminer la position de l'image et le grandissement.
\item Construire l'image et expliquer pourquoi on ne peut pas la recueillir directement sur un écran à sa position.
\end{questions}
\ex{E12}{Une divergente en autonomie}{Niveau 2}{15 min}
Un objet réel est à $36{,}0\ \mathrm{cm}$ d'une lentille de focale $-12{,}0\ \mathrm{cm}$.
\begin{questions}
\item Calculer position et grandissement de l'image ; vérifier les signes sur une construction.
\item Montrer plus généralement que l'image d'un objet réel par une lentille divergente est virtuelle, droite et réduite.
\end{questions}
\begin{docu}{Grille commune de correction}
\begin{tabularx}{\linewidth}{@{}Y Y Y@{}}\toprule
Position & Orientation & Taille\\\midrule
$q>0$ : réelle & $\gamma<0$ : renversée & $|\gamma|>1$ : agrandie\\
$q<0$ : virtuelle & $\gamma>0$ : droite & $|\gamma|<1$ : réduite\\\bottomrule
\end{tabularx}
Ces critères supposent l'orientation d'axe et le sens de propagation fixés dans le cours. La construction fournit un contrôle indépendant du calcul.
\end{docu}

\Page{Exercices\quad /\quad Entraînement}{Foyers et projection}{td5}
\ex{E13}{Une étoile hors de l'axe}{Niveau 2}{12 min}
Un point à l'infini envoie sur une lentille convergente de focale $200\ \mathrm{mm}$ un faisceau dont la direction fait un angle $+0{,}500^\circ$ avec l'axe, orienté vers le haut.
\begin{questions}
\item Construire son image et préciser le plan qui la contient.
\item Calculer sa hauteur algébrique par rapport à l'axe dans l'approximation de Gauss.
\item Où placer un point source pour obtenir, en sens inverse, un faisceau parallèle de même inclinaison ?
\end{questions}
\ex{E14}{Un objet très proche du foyer}{Niveau 2}{12 min}
Une lentille convergente a $f'=5{,}00\ \mathrm{mm}$. Un objet réel est situé $1{,}00\ \mathrm{mm}$ avant son foyer objet.
\begin{questions}
\item Donner $\overline{FA}$ et utiliser Newton pour déterminer $\overline{F'A'}$.
\item En déduire $\oap$ et $\gamma$, puis contrôler par Descartes.
\item Que devient l'image si l'objet se rapproche du foyer par la gauche ?
\end{questions}
\ex{E15}{Deux positions ou aucune}{Niveau 2}{20 min}
Un objet et un écran sont séparés de $90{,}0\ \mathrm{cm}$. On dispose d'une lentille de focale $20{,}0\ \mathrm{cm}$.
\begin{questions}
\item Sans résoudre l'équation, justifier que deux positions sont possibles.
\item Déterminer les deux distances objet-lentille et les grandissements.
\item Reprendre la conclusion si la distance objet-écran vaut $75{,}0\ \mathrm{cm}$, puis $80{,}0\ \mathrm{cm}$.
\end{questions}
\begin{methode}{Une formule n'est pas toujours le plus court chemin}
Dans E14, l'origine au foyer évite une soustraction peu lisible. Dans E15, le discriminant renseigne sur l'existence avant de donner les valeurs.
\end{methode}

\Page{Exercices\quad /\quad Applications}{Voir et photographier}{td6}
\ex{E16}{L'accommodation chiffrée}{Niveau 1}{12 min}
Un œil réduit emmétrope possède une distance lentille-rétine $\ell=17{,}0\ \mathrm{mm}$ et voit net de $25{,}0\ \mathrm{cm}$ à l'infini.
\begin{questions}
\item Donner les positions du punctum remotum et du punctum proximum.
\item Calculer les vergences extrêmes et leur différence.
\item Pourquoi un modèle à focale constante et rétine mobile ne décrit-il pas l'accommodation de cet œil ?
\end{questions}
\ex{E17}{Distinguer deux détails}{Niveau 1}{10 min}
Deux traits espacés de $2{,}00\ \mathrm{cm}$ sont observés à $100\ \mathrm{m}$. La résolution angulaire de l'observateur est de $1{,}00$ minute d'arc.
\begin{questions}
\item Convertir cette limite en radians.
\item Les traits sont-ils résolus avec le critère $\theta\geq\theta_{\min}$ ?
\item À quelle distance maximale pourraient-ils être distingués ?
\end{questions}
\ex{E18}{Un portrait qui ne tient pas}{Niveau 2}{18 min}
Un appareil est modélisé par une lentille de focale $50{,}0\ \mathrm{mm}$ et un capteur de hauteur $24{,}0\ \mathrm{mm}$. Une personne de hauteur $1{,}80\ \mathrm{m}$ se trouve à $2{,}00\ \mathrm{m}$ de la lentille.
\begin{questions}
\item Calculer la distance de mise au point et son écart à la position pour un objet à l'infini.
\item Calculer la taille de l'image et déterminer si la personne tient entièrement dans le cadre.
\item À quelle distance minimale placer cette personne pour que son image ait une hauteur inférieure ou égale à celle du capteur ?
\end{questions}
\source{Pour E18, les distances de l'objet et du capteur se mesurent à partir du centre optique du modèle. On ne remplace pas systématiquement $q$ par $f'$ pour un objet à distance finie.}

\Page{Exercices\quad /\quad Applications}{Ouverture et fibre}{td7}
\ex{E19}{Exposition et mouvement}{Niveau 2}{15 min}
Un réglage $N=2{,}8$ et $T=1/500\ \mathrm{s}$ donne une exposition correcte. L'image d'un détail mobile se déplace sur le capteur à $10{,}0\ \mathrm{mm\,s^{-1}}$ ; le pas de pixel est $4{,}00\ \mu\mathrm{m}$.
\begin{questions}
\item Quelle durée conserve l'exposition lorsque l'on ferme à $N=5{,}6$ ?
\item Calculer le flou de mouvement, en micromètres puis en pixels, pour chaque réglage.
\item Quel compromis faut-il expliquer avant de choisir le réglage ?
\end{questions}
\ex{E20}{Un point encore acceptable}{Niveau 2}{18 min}
Un objectif de focale $50{,}0\ \mathrm{mm}$ est réglé sur $p_0=5{,}00\ \mathrm{m}$, à $N=8{,}00$. Un point de l'axe est à $p=4{,}00\ \mathrm{m}$. Le diamètre admissible de flou est $\varepsilon=30{,}0\ \mu\mathrm{m}$.
\fig{ex-profondeur}{Compléter le cône issu de $A$ après la lentille, puis sa trace sur le capteur ; schéma sans échelle.}
\begin{questions}
\item Calculer le diamètre de l'ouverture, le plan du capteur $q_0$ et le plan image idéal $q$ du point.
\item Construire les deux rayons extrêmes et calculer $b=a|1-q_0/q|$.
\item Le point est-il acceptable ? Que devient $b$ si l'on ferme à $N=16$, à réglage de mise au point inchangé ?
\end{questions}
\ex{E21}{Injecter dans une fibre}{Niveau 2}{15 min}
Une fibre possède $n_1=1{,}48$, $n_2=1{,}46$ et une face d'entrée normale à son axe.
\begin{questions}
\item Calculer son ouverture numérique et son angle d'acceptance dans l'air.
\item Un rayon arrivant à $10{,}0^\circ$ de l'axe est-il guidé dans le modèle méridien ? Justifier à la paroi.
\item Que devient l'angle d'acceptance si le milieu d'injection est l'eau d'indice $1{,}33$ ?
\end{questions}

\Page{Exercices\quad /\quad Approfondissement}{Assembler et vérifier}{td8}
\ex{E22}{Une impulsion qui s'étale}{Niveau 2}{18 min}
Une fibre rectiligne de $1{,}00\ \mathrm{km}$ a les indices $n_1=1{,}48$ et $n_2=1{,}46$. On néglige toute dispersion matérielle.
\begin{questions}
\item Établir $t(\theta)$, puis calculer $t_{\min}$ et $\Delta t$.
\item Quel débit maximal estime-t-on si la période binaire doit être au moins égale à $2\Delta t$ ?
\item Distinguer le retard de propagation et l'étalement de l'impulsion.
\end{questions}
\ex{E23}{Deux lentilles successives}{Niveau 3}{25 min}
Un objet est à $30{,}0\ \mathrm{cm}$ à gauche d'une lentille $L_1$ de focale $+10{,}0\ \mathrm{cm}$. Une lentille $L_2$ de focale $-10{,}0\ \mathrm{cm}$ est à $20{,}0\ \mathrm{cm}$ à droite de $L_1$.
\fig{ex-association}{Les deux lentilles et leurs foyers : ajouter l'image intermédiaire puis l'image finale.}
\begin{questions}
\item Trouver l'image intermédiaire $A_1B_1$ donnée par $L_1$.
\item Donner $\overline{O_2A_1}$ et déterminer l'image finale, sa nature et sa position par rapport à $O_1$.
\item Calculer le grandissement total et expliquer pourquoi les deux origines ne doivent pas être confondues.
\end{questions}
\ex{E24}{Un résultat de mesure défendable}{Niveau 3}{20 min}
Une lentille donne une image nette pour $p=30{,}0\ \mathrm{cm}$ et $q=60{,}0\ \mathrm{cm}$. Ces deux distances sont ici mesurées indépendamment, avec $u(p)=u(q)=0{,}10\ \mathrm{cm}$.
\begin{questions}
\item Calculer la focale, le grandissement et la vergence.
\item Propager les incertitudes à l'aide des dérivées partielles de $f'=pq/(p+q)$.
\item Un résultat indépendant vaut $f'_r=20{,}10\ \mathrm{cm}$ avec $u(f'_r)=0{,}08\ \mathrm{cm}$. Calculer l'écart normalisé et conclure avec le repère $z\leq2$.
\end{questions}
